\begin{figure}
    \centering
    \begin{tikzpicture}[
        scale=0.85,
        latticepoint/.style={circle, fill=black!55,
                             inner sep=1.2pt},
        basis/.style={-{Stealth[length=2.5mm]},
                      line width=1.2pt}
    ]

    % Short basis
    \begin{scope}
        % Lattice points
        \foreach \i in {-3,...,3} {
            \foreach \j in {-3,...,3} {
                \node[latticepoint] at
                    ({\i+0.35*\j},{\j}) {};
            }
        }

        % Basis vectors
        \draw[basis, teal!80!black]
            (0,0) -- (1,0);
        \draw[basis, teal!80!black]
            (0,0) -- (0.35,1);

        \node[below] at (0,-3.65)
            {\textbf{Short basis}};
    \end{scope}

    % Long basis
    \begin{scope}[xshift=8cm]
        % Identical lattice points
        \foreach \i in {-3,...,3} {
            \foreach \j in {-3,...,3} {
                \node[latticepoint] at
                    ({\i+0.35*\j},{\j}) {};
            }
        }

        % Long basis vectors
        % b1 = 3v1 + 2v2
        % b2 = 2v1 + v2
        \draw[basis, orange!85!black]
            (0,0) -- (3.7,2);
        \draw[basis, orange!85!black]
            (0,0) -- (2.35,1);

        \node[below] at (0,-3.65)
            {\textbf{Long basis}};
    \end{scope}

    \end{tikzpicture}
    \caption{Two different bases of the same lattice.
    A short  \textit{easy} basis (left) and a long \textit{hard}
basis (right).}
    \label{fig:lattice-bases}
\end{figure}